\documentclass[aspectratio=169,12pt]{beamer}

\usetheme{default}
\usecolortheme{default}
\setbeamertemplate{navigation symbols}{}

\setbeamercolor{background canvas}{bg=white}
\setbeamercolor{normal text}{fg=black,bg=white}
\setbeamercolor{frametitle}{fg=black,bg=white}
\setbeamercolor{title}{fg=black}
\setbeamercolor{subtitle}{fg=black}
\setbeamercolor{author}{fg=black}
\setbeamercolor{institute}{fg=black}
\setbeamercolor{date}{fg=black}
\setbeamercolor{item}{fg=black}
\setbeamercolor{subitem}{fg=black}
\setbeamercolor{enumerate item}{fg=black}
\setbeamertemplate{blocks}[default]
\setbeamercolor{block title}{fg=black,bg=white}
\setbeamercolor{block body}{fg=black,bg=white}

\setbeamertemplate{frametitle}{
  \vspace{0.4em}
  \insertframetitle\par
  \vspace{0.15em}
  {\color{black!55}\rule{\textwidth}{0.4pt}}
  \vspace{0.2em}
}

\setbeamertemplate{footline}{
  \hfill
  {\color{black!55}\insertframenumber/\inserttotalframenumber}
  \hspace{1.2em}
  \vspace{0.6em}
}

\usepackage[utf8]{inputenc}
\usepackage[T2A]{fontenc}
\usepackage[russian]{babel}
\usepackage{graphicx}
\usepackage{hyperref}
\hypersetup{colorlinks=false,pdfborder={0 0 0}}

\graphicspath{{../../report/images/}{../images/}}

\newcommand{\screenshotslide}[3]{%
  \begin{frame}{#1}
    \begin{center}
      \includegraphics[width=0.90\linewidth,height=0.64\textheight,keepaspectratio]{#2}
      \vspace{0.35em}
      \parbox{0.94\linewidth}{\centering\small #3}
    \end{center}
  \end{frame}
}

\title{Разработка интерактивной аналитической панели\\для визуализации и анализа данных\\рынка недвижимости}
\subtitle{на основе данных платформы Яндекс.Недвижимость}
\author{Перфилов Арсений Сергеевич}
\institute{341 группа}
\date{2026}

\begin{document}

% ===== I. ВВЕДЕНИЕ =====

\begin{frame}[plain]
  \vfill
  {\Large\bfseries\inserttitle\par}
  \vspace{0.8em}
  {\normalsize\insertsubtitle\par}
  \vspace{2.2em}
  {\insertauthor\par}
  \vspace{0.4em}
  {\small\insertinstitute\par}
  \vspace{0.4em}
  {\small\insertdate\par}
  \vfill
\end{frame}

\begin{frame}{Актуальность}
  \begin{itemize}
    \item Тысячи объявлений на Яндекс.Недвижимость --- мало удобных инструментов для анализа.
    \item Нельзя быстро сравнить сегменты рынка и посмотреть выборку на карте.
    \item Нет публичного API; данные динамические, работает антибот-защита.
  \end{itemize}
\end{frame}

\begin{frame}{Цель и задачи}
  \textbf{Цель} --- система полного цикла: сбор $\rightarrow$ хранение $\rightarrow$ интерактивный анализ.

  \medskip
  \textbf{Задачи:}
  \begin{enumerate}
    \item Парсер (Python, Selenium): ссылки $\rightarrow$ карточки $\rightarrow$ CSV.
    \item Сервер (Express, SQLite): датасеты по городам, REST API.
    \item Дашборд (React): фильтры, KPI, графики, карта, таблица.
    \item Роли пользователей; проверка на Москве и Саратове.
  \end{enumerate}
\end{frame}

\begin{frame}{Архитектура системы}
  \begin{center}
    \setlength{\tabcolsep}{8pt}
    \begin{tabular}{c}
      \textbf{1. Парсер} --- Python, Selenium, ETL \\[0.3em]
      $\downarrow$ CSV \\[0.3em]
      \textbf{2. Сервер} --- Express, SQLite, REST API \\[0.3em]
      $\downarrow$ JSON \\[0.3em]
      \textbf{3. Клиент} --- React, Recharts, Яндекс.Карты \\
    \end{tabular}
  \end{center}
  \vfill
  \small Общие типы и фильтры --- каталог \texttt{shared/}. Далее --- реализация по этим трём слоям.
\end{frame}

% ===== II. СБОР И ХРАНЕНИЕ ДАННЫХ =====

\screenshotslide{Сбор данных: парсер}{admin-parser-running.png}{Этап «Сбор ссылок»: прогресс, возобновление после остановки, обработка капчи}

\screenshotslide{Хранение данных: сервер}{admin-server.png}{Загрузка CSV, ingest в датасет города, история обновлений}

% ===== III. АНАЛИЗ В ДАШБОРДЕ (один сценарий) =====

\screenshotslide{Дашборд: исходная выборка}{dash-main-light-moscow.png}{Москва, 12\,549 объявлений --- фильтры, KPI и графики на полном массиве}

\screenshotslide{Шаг 1 --- фильтрация}{filters-after-light.png}{Срез: 3 комнаты, 120--190~м$^2$, 30--45~млн~руб.\\Осталось 22 объекта из 12\,549}

\screenshotslide{Шаг 2 --- KPI и сегменты}{charts-after-filter-22.png}{Пересчёт показателей; сравнение: 7 новостроек и 15 вторичных}

\screenshotslide{Шаг 3 --- карта и таблица}{map-table-after-filter.png}{19 точек на карте, 22 строки в таблице; данные согласованы с API}

% ===== IV. РОЛИ ПОЛЬЗОВАТЕЛЕЙ =====

\begin{frame}{Роли пользователей}
  \begin{itemize}
    \item \textbf{observer} --- просмотр KPI, графиков и карты; упрощённый список.
    \item \textbf{analyst} --- полный дашборд: фильтры, экспорт, конструктор графиков.
    \item \textbf{admin} --- датасеты, ingest и запуск парсера (слайды 5--6).
    \item \textbf{manager} --- обзор рынков по городам и управление пользователями.
  \end{itemize}
  \vfill
  \small Ниже --- примеры интерфейсов для manager и observer.
\end{frame}

\screenshotslide{Роль manager}{manager-markets-dark.png}{Сводка по городам: число объявлений и медианная цена}

\screenshotslide{Роль observer}{observer-saratov-light.png}{Саратов, 75 объявлений --- облегчённый интерфейс без полной таблицы}

% ===== V. ИТОГИ =====

\begin{frame}{Результаты проверки}
  \begin{itemize}
    \item Москва --- \textbf{12\,549} объявлений; Саратов --- \textbf{75}.
    \item Пройден сценарий фильтрации (22 объекта) и все роли.
    \item Развёртывание на VPS: nginx, systemd.
    \item Дашборд: \texttt{http://62.217.176.204/}
  \end{itemize}
\end{frame}

\begin{frame}{Итоги}
  Разработан программный комплекс полного цикла для анализа рынка жилой недвижимости на данных Яндекс.Недвижимость: парсер, сервер, интерактивный дашборд с разграничением прав.

  \vfill
  \centering
  \href{http://62.217.176.204/}{http://62.217.176.204/}
\end{frame}

\end{document}
